\section{Method}
\label{sec:method}

\subsection{Task Formulation}
\label{sec:task}

We study the task of translating natural language (NL) premises and hypotheses into formal logic representations verifiable by an external solver.
Given NL premises $P$ and hypothesis $H$, a formalizer $\pi_\theta$ generates a formal representation $x$ comprising $n$ components $\{c_1, \ldots, c_n\}$---axioms, rules, and a conclusion---submitted to solver $\mathcal{S}$, which returns a binary outcome: \texttt{proved} or \texttt{fail}.

We consider two target formalisms:
(1)~\textbf{First-order logic (FOL)}, where components are Prover9-format axioms and a goal clause, verified by the Prover9 theorem prover~\cite{mccune2010prover9};
(2)~\textbf{Prolog}, where components are Horn clauses, facts, and a query, verified by SWI-Prolog.
Training data is drawn from FOLIO~\cite{folio}, a corpus of 1{,}204 human-annotated FOL examples.
We use a fixed split of 500 training and 100 held-out evaluation examples across all experiments.

\subsection{GRPO Training}
\label{sec:grpo}

We train the formalizer with Group Relative Policy Optimization~\citep[GRPO;][]{deepseekmath}.
For each training problem, the policy $\pi_\theta$ generates $K{=}8$ candidate formalizations.
Each candidate is submitted to the solver, yielding a binary execution reward:
\begin{equation}
  r_{\text{exec}}(x) =
  \begin{cases}
    1 & \text{if } \mathcal{S}(x) = \texttt{proved} \\
    0 & \text{otherwise}
  \end{cases}
  \label{eq:r_exec}
\end{equation}
GRPO normalizes rewards within each group of $K$ completions to estimate advantages, then updates $\pi_\theta$ via a clipped policy gradient objective.
The base model is Qwen3-4B; all runs use learning rate $10^{-6}$ and train for 500 gradient steps unless otherwise noted.

Under execution-only reward, the optimizer receives no signal about \emph{why} a proof succeeds---whether the formalization captures genuine logical derivation or exploits a degenerate shortcut.
This opacity motivates a structural diagnostic.

\subsection{Solver Trace Ablation}
\label{sec:sta}

Solver Trace Ablation (STA) tests whether each component of a proved formalization is \emph{necessary} for the proof.
The protocol applies leave-one-out (LOO) ablation---a standard technique in mutation testing~\cite{demillo1978, jia-harman-survey} and feature attribution---to solver inputs.
We claim no novelty in the ablation mechanism itself; the contribution lies in applying it as a diagnostic for reward hacking in neurosymbolic GRPO training.

\begin{algorithm}[t]
\caption{Solver Trace Ablation (STA)}
\label{alg:sta}
\begin{algorithmic}[1]
\Require Proved formalization $x = \{c_1, \ldots, c_n\}$, solver $\mathcal{S}$
\Ensure $\text{ISR}(x) \in [0, 1]$
\State $\textit{necessary} \gets 0$
\For{$i = 1$ to $n$}
    \State $x_{\setminus i} \gets x \setminus \{c_i\}$
    \If{$\mathcal{S}(x_{\setminus i}) = \texttt{fail}$}
        \State $\textit{necessary} \gets \textit{necessary} + 1$
    \EndIf
\EndFor
\State \Return $\textit{necessary} \;/\; n$
\end{algorithmic}
\end{algorithm}

For a proved formalization $x = \{c_1, \ldots, c_n\}$, STA removes each component in turn and re-runs the solver (Algorithm~\ref{alg:sta}).
The protocol requires $n$ solver calls per proved sample.
For unproved samples ($\mathcal{S}(x) = \texttt{fail}$), ISR is undefined: there is no successful proof to ablate.
Throughout the paper, we report \emph{proved-only} ISR (denoted ISR$^*$), computed exclusively over samples the solver proves, to avoid conflating structural quality with prove rate.
Figure~\ref{fig:sta-overview} illustrates the full protocol.

\subsection{Intensional Structural Reward}
\label{sec:isr}

The \textbf{Intensional Structural Reward} (ISR) quantifies the fraction of components whose removal causes proof failure:
\begin{equation}
  \text{ISR}(x) = \frac{1}{n} \sum_{i=1}^{n} \mathbf{1}\!\left[\mathcal{S}(x_{\setminus i}) = \texttt{fail}\right]
  \label{eq:isr}
\end{equation}

ISR is bounded in $[0, 1]$.
$\text{ISR} = 1$ indicates every component is individually necessary---removing any single one breaks the proof.
$\text{ISR} = 0$ indicates no component is individually necessary, signaling redundancy or degeneracy.

ISR measures \emph{syntactic component necessity}: whether each piece is required for the proof to succeed.
It does not measure semantic faithfulness to the NL input.
A formalization can achieve $\text{ISR} = 1$ while being semantically unrelated to the original statement---a single premise identical to the conclusion trivially satisfies necessity.

\paragraph{Premise count bias.}
ISR is mathematically sensitive to component count $n$.
When $n = 2$ (one premise, one conclusion), removing the sole premise breaks any non-trivial proof, yielding $\text{ISR} \geq 0.5$ regardless of proof quality.
Conversely, formalizations with many components tend toward lower ISR as redundancy accumulates.
This is an inherent property of LOO-based metrics, not specific to our setting.
When comparing ISR across conditions that differ in mean premise count, we control for this confound via premise-count-matched analysis (\S\ref{sec:experiments}).

\subsection{Reward Formulations}
\label{sec:reward}

We study six reward conditions to isolate the effect of ISR as a training signal (Table~\ref{tab:reward}).

\begin{table}[t]
\centering
\small
\begin{tabular}{@{}llc@{}}
\toprule
\textbf{Condition} & \textbf{Reward formula} & $\beta$ \\
\midrule
Exec-only (baseline) & $r_{\text{exec}}$ & --- \\
ISR low-dose & $r_{\text{exec}} \cdot (1 + \beta \cdot \text{ISR})$ & 0.1 \\
ISR high-dose & $r_{\text{exec}} \cdot (1 + \beta \cdot \text{ISR})$ & 0.3 \\
Random low-dose & $r_{\text{exec}} \cdot (1 + \beta \cdot u)$ & 0.1 \\
Random high-dose & $r_{\text{exec}} \cdot (1 + \beta \cdot u)$ & 0.3 \\
SFT baseline & Supervised (NLL on GT) & --- \\
\bottomrule
\end{tabular}
\caption{Reward conditions. $u \sim \text{Uniform}(0,1)$. All GRPO conditions use Qwen3-4B, $K{=}8$, lr$=10^{-6}$, 500 steps.}
\label{tab:reward}
\end{table}

\paragraph{Additive ISR bonus.}
ISR augments the execution reward as a multiplicative bonus, active only for proved samples:
\begin{equation}
  r = r_{\text{exec}} \cdot (1 + \beta \cdot \text{ISR}(x))
  \label{eq:additive_isr}
\end{equation}
Proved samples receive reward in $[1, 1{+}\beta]$; unproved samples receive $0$.
We evaluate $\beta \in \{0.1, 0.3\}$ to test dose-response dynamics.
An earlier convex formulation ($r = \alpha \cdot r_{\text{exec}} + \beta \cdot r_{\text{exec}} \cdot \text{ISR}$, $\alpha{+}\beta{=}1$) was discarded because $\text{ISR} < 1$ always penalized relative to pure execution reward, collapsing prove rate.

\paragraph{Random additive control.}
To disentangle ISR's structural signal from the effect of reward variance, we replace ISR with a uniform random variable:
\begin{equation}
  r = r_{\text{exec}} \cdot (1 + \beta \cdot u), \quad u \sim \text{Uniform}(0, 1)
  \label{eq:additive_random}
\end{equation}
with identical $\beta$ values.
If ISR provides genuine structural guidance beyond stochastic perturbation, it should outperform this control.

\paragraph{SFT baseline.}
Standard supervised fine-tuning on the same 500 training examples using ground-truth FOL annotations as targets.
No reinforcement learning or reward shaping.
This control isolates whether structural degeneration stems from RL optimization or from data exposure alone.

All GRPO conditions share identical hyperparameters and are evaluated across seeds $\in \{42, 43, 44\}$.

\subsection{Hacking Taxonomy}
\label{sec:taxonomy}

We identify seven distinct modes of structural degeneration across reward conditions and logic types.
Definitions appear below; quantitative prevalence analysis is deferred to~\S\ref{sec:experiments}.

\textbf{H1: Conclusion embedding.}
At least one premise is identical to the conclusion after normalization, making the proof trivially $p \rightarrow p$.
This is the dominant mode across all GRPO conditions (52--93\% of proved samples).

\textbf{H7: Conclusion duplication with ISR inflation.}
The formalization contains a single premise identical to the conclusion ($n{=}2$).
Removing the premise breaks the proof, so $\text{ISR} \geq 0.5$ by construction.
This mode is exclusive to ISR-reward conditions (64.5\% vs.\ 0\% under random control across three seeds) and represents ISR-specific gaming.

\textbf{Tautological collapse.}
Degeneration to a minimal two-component structure that is trivially provable, arising under exec-only conditions without ISR inflation.

\textbf{Repetition collapse.}
Identical premises generated repeatedly, inflating component count without adding logical content.

\textbf{Complexity without utility.}
Syntactically complex but logically vacuous premises that do not contribute to the proof.

\textbf{Contradiction explosion.}
Contradictory premises (e.g., $p(a) \wedge \neg p(a)$) from which any conclusion follows \emph{ex falso quodlibet}.
Observed at extended training (1000 steps), with mean component count rising from 6.6 to 18.6.

\textbf{Query-as-fact gaming} (Prolog-specific).
The model converts the NL question into a Prolog atom asserted simultaneously as fact and query, making the proof trivially true.
This mode evades text-matching classifiers designed to detect conclusion embedding; detection requires functional analysis of the STA necessity vector ($n_{\text{necessary}} = 1$ for 95/97 proved samples).

Detection combines two complementary methods: normalized text matching (for H1, H7) and functional analysis of the necessity vector (for modes that bypass text-based classifiers, such as query-as-fact gaming).
We report both strict (text-match) and functional ($n_{\text{necessary}}$-based) classifications throughout.
